% ECONOMICS_PROBLEM: {"id": "C18-215", "title": "Partial identification under model misspecification", "jel_code": "C18", "source_id": "OP-0110"}
% !TeX program = xelatex
\documentclass[11pt,a4paper]{article}
\usepackage[margin=23mm]{geometry}
\usepackage{fontspec}
\setmainfont{Latin Modern Roman}
\usepackage{amsmath,amssymb,mathtools}
\usepackage{xcolor,enumitem,booktabs,longtable,array,xurl,hyperref}
\definecolor{muted}{HTML}{53636C}
\hypersetup{colorlinks=true,urlcolor=blue,linkcolor=blue}
\setlength{\parindent}{0pt}
\setlength{\parskip}{5pt}
\newcommand{\R}{\mathbb R}
\newcommand{\E}{\mathbb E}
\newcommand{\N}{\mathbb N}
\newcommand{\ind}{\mathbf 1}
\newcommand{\Var}{\operatorname{Var}}
\newcommand{\Cov}{\operatorname{Cov}}
\newcommand{\supp}{\operatorname{supp}}
\DeclareMathOperator*{\argmin}{arg\,min}
\DeclareMathOperator*{\argmax}{arg\,max}
\newcommand{\entrymeta}[3]{{\sffamily\small\color{muted}\textbf{\texttt{#1}}\quad|\quad #2\par #3\par}}
\newcommand{\Problemref}[1]{\texttt{#1}}
\newcommand{\JELref}[2]{\texttt{#2} (JEL \texttt{#1})}
\begin{document}
\section*{Partial identification under model misspecification}
\textbf{Problem ID:} \texttt{C18-215}\quad \textbf{JEL:} \texttt{C18}\par
% BEGIN ECONOMICS STATEMENT
\label{problem:OP-0110}
\label{atlas:prob:E10}

\noindent\textbf{Original identifier:} E10 (atlas item 110). \textbf{Field:} Econometrics and causal inference. \textbf{Classification:} editorial research family with established restricted solutions; current open status not certified.

\subsection*{Definitions and assumptions}

Let $P_0$ be the observable-data law and let a structural model $m\in\mathcal M$ induce law $P_m$ and policy value $V(a,m)$ for decision $a\in\mathcal A$. Define a supplied statistical distance $d$ and tolerance $\varepsilon\ge0$. The observationally admissible class is $\mathcal M(P_0,\varepsilon)=\{m\in\mathcal M:d(P_m,P_0)\le\varepsilon\}$. For a scalar structural target $\theta(m)$, its identified or sensitivity set is $\Theta(P_0,\varepsilon)=\{\theta(m):m\in\mathcal M(P_0,\varepsilon)\}$. Specify compactness, continuity, bounded welfare and measurable selections if existence of optimizers is asserted.

\subsection*{Formal research question}

Construct inference and decision rules robust to a declared degree of observable misspecification and structural ambiguity. Given a confidence region $\mathcal C_n$ for $P_0$, define the model envelope $\widehat{\mathcal M}_n=\bigcup_{P\in\mathcal C_n}\mathcal M(P,\varepsilon_n)$. Seek coverage of $\Theta(P_0,\varepsilon_n)$ uniformly over a specified sampling class, including weakly binding moments and empty baseline models. When $\widehat{\mathcal M}_n$ is nonempty, a robust decision objective is
\[
\widehat a_n\in\arg\min_{a\in\mathcal A}
\sup_{m\in\widehat{\mathcal M}_n}
\left[\sup_{b\in\mathcal A}V(b,m)-V(a,m)\right].
\]
Characterize regret relative to decisions using the population admissible class, and derive rates separating sampling error, approximation tolerance and unavoidable ambiguity. State how $\varepsilon_n$ is fixed or selected without treating a data-dependent fit threshold as an external truth.

\subsection*{Interpretation and scope}

Partial identification within a correct model and robustness to a false model are distinct. A model may fit observables yet predict counterfactual policy outcomes incorrectly; the distance restriction alone does not validate counterfactual invariance. Empty identified sets are possible evidence of refutation, not estimates of zero effects. A tolerance can prevent emptiness but may admit substantively implausible models, so restrictions need economic content. Local sensitivity derivatives summarize small perturbations and need not bound large departures. A minimax-regret objective is one defensible decision criterion; Bayesian averaging and maximin welfare embody other attitudes to ambiguity and must not be silently substituted. If the model envelope is empty, report refutation or use a prespecified enlargement; the displayed supremum-based decision problem is not applied to an empty set.

\subsection*{Source audit and status}

The slash-year Manski item is split into the verified 1990 article [S1] and 2003 book [S2]. Tamer [S3] reviews partial identification; Andrews--Gentzkow--Shapiro [S4] studies local moment sensitivity. The input's author-homepage link is not an exact source for its decision-theory claim. The robust envelope-and-regret formulation is editorial and extends the cited backgrounds; these works do not certify this precise research question as open in 2026.

\subsection*{References}

\begin{enumerate}[label={[\textnormal{S}\arabic*]},leftmargin=*]

\item Charles F. Manski (1990). \emph{Nonparametric Bounds on Treatment Effects}. American Economic Review 80(2), 319–323. \url{https://www.jstor.org/stable/2006592}. \textbf{Locator:} Original 1990 article scan, sections I–II; canonical JSTOR article landing page. \textbf{Support and limits:} Author-authored article scan verifies original bounds paper; not a misspecification-robust theorem.

\item Charles F. Manski (2003). \emph{Partial Identification of Probability Distributions}. Springer. \url{https://doi.org/10.1007/b97478}. \textbf{Locator:} Publisher or institutional abstract and bibliographic record. \textbf{Support and limits:} Book develops partial-identification analysis; this is a precise completion of the slash-year citation.

\item Elie Tamer (2010). \emph{Partial Identification in Econometrics}. Annual Review of Economics 2, 167–195. \url{https://doi.org/10.1146/annurev.economics.050708.143401}. \textbf{Locator:} Publisher or institutional abstract and bibliographic record. \textbf{Support and limits:} Survey of identified sets and inference; partial identification does not itself cure a false model.

\item Isaiah Andrews, Matthew Gentzkow and Jesse M. Shapiro (2017). \emph{Measuring the Sensitivity of Parameter Estimates to Estimation Moments}. Quarterly Journal of Economics 132(4), 1553–1592. \url{https://academic.oup.com/qje/article-abstract/132/4/1553/3861634}. \textbf{Locator:} Publisher or institutional abstract and bibliographic record. \textbf{Support and limits:} Local sensitivity of estimators to moments; not global robust decision theory.

\end{enumerate}

\subsection*{Common conventions: Source mathematical conventions}

Unless an entry states otherwise, agent and firm sets are finite. State and action spaces are measurable, strategies and outcome maps are measurable, probability laws are specified, and expectations exist for the targets under discussion. Infinite-horizon discount factors lie in $(0,1)$ and discounted objectives are finite. Norms, tolerances and complexity input sizes must be specified for computational comparisons. Constants or primitives that appear locally are inputs, not empirical facts asserted by this catalogue.

\textbf{Identification.} A statistical model $m\in\mathcal M$ implies an observable law $P_m$ and target $\theta(m)$. At law $P$, its identified set is
\[
\Theta(P;\mathcal M)=\{\theta(m):m\in\mathcal M,\ P_m=P\}.
\]
Point identification holds when this set is a singleton, or equivalently when $P_{m_1}=P_{m_2}$ implies $\theta(m_1)=\theta(m_2)$ for all compatible models. Sharp bounds mean bounds attained, or approached when the boundary is not attained, by compatible models. Writing this set is a definition, not a solution: the substantive task is to establish informative restrictions and derive or estimate the set. An unrestricted-confounding model may give only trivial bounds.

\textbf{Inference.} Identification, estimability and valid uncertainty quantification are separate questions. Fix $\alpha\in(0,1)$. A confidence procedure $C_n$ over a declared law class $\mathcal P$ has asymptotic uniform coverage at level $1-\alpha$ when
\[
\liminf_{n\to\infty}\inf_{P\in\mathcal P}\Pr_P\{\theta(P)\in C_n\}\geq1-\alpha.
\]
For set-identified targets the coverage event must be defined separately, for example coverage of the full identified set. If an entry uses identification-indexed models instead of uniquely indexed laws, the same coverage convention is applied over those models. Pointwise limits do not establish uniform validity, and asymptotic validity does not guarantee finite-sample precision. Sample sizes, dependence structures and weak-identification sequences are part of each application.

\textbf{Counterfactuals.} $Y_i(a)$ is a potential outcome under a specified intervention, including its timing, implementation and versions. It is not $Y_i$ conditional on voluntarily choosing $a$. A full intervention vector permits interference; shorthand scalar treatment notation does not silently impose no interference. Assignment, selection, attrition, anticipation and measurement must be specified. A claim about an instrument's local effect is not automatically a population policy effect.

\textbf{Economic equilibrium.} A counterfactual model includes best responses, constraints, feasibility, market clearing, state transitions and expectations consistent with the stated information. When several equilibria exist, either a declared selector is an input or the target is set-valued. Comparisons across policy regimes must use the stated selection convention. Reduced-form relationships are not presumed invariant after an equilibrium-changing intervention.

\textbf{Optimization and welfare.} Optimization is over the stated feasible set. If attainment is not established, $\sup$ or $\inf$ and $\varepsilon$-optimal choices replace an asserted maximizer or minimizer. For example, with value $V=\sup_{a\in\mathcal A}W(a)<\infty$, an $\varepsilon$-optimal set is $\{a\in\mathcal A:W(a)\geq V-\varepsilon\}$ for $\varepsilon>0$. Benefits, costs, risk and health or environmental outcomes can be aggregated only using declared units and conversion weights. Interpersonal utility weights, the treatment of fiscal transfers and foreign profits, discounting, and the behavioral welfare criterion are explicit inputs. A comparative-static derivative additionally requires existence and appropriate differentiability of the selected counterfactual.

\textbf{Sources.} Local labels [S1], [S2], and so forth restart in every question. Locators identify the relevant abstract, model, section or result at the level actually checked; a bibliographic or abstract-level verification is not represented as inspection of an entire proof. Working papers are distinguished from later publications and changing versions. Source summaries are paraphrased. All URL checks and editorial formulations refer to the audit date stated above.
% END ECONOMICS STATEMENT
\end{document}
